\begin{rSection}{Experience}

\begin{rSubsection}{Keysight Technologies}{May 2026 -- Dec 2026}{Machine Learning Security Intern, Applied ML \& Fault Injection}{Santa Rosa, CA}
\item \textbf{Scalable fault-injection tooling:} built model-guided optimization and adaptive-learning workflows that automatically tune fault-injection parameters, improving vulnerability-detection coverage over baseline sweep-based flows (GlitchLab, arXiv 2026).
\item \textbf{Experiment and telemetry design:} designing experiments, telemetry instrumentation, scalable hyperparameter search, and reliability metrics; delivering a prototype demonstration to the security-test team.
\end{rSubsection}

\begin{rSubsection}{The University of Kansas}{Jan 2022 -- Present}{Graduate Research Assistant, Hardware Security (Advisor: Dr.\ Tamzidul Hoque)}{Lawrence, KS}
\item \textbf{Side-channel and fault attacks on cryptography:} demonstrated power side-channel and fault attacks on the ASCON lightweight-cryptography standard and evaluated countermeasures.
\item \textbf{Pre-deployment security testing:} proposed a gray-box evaluation methodology that uses power and EM side-channel measurements to flag anomalous behavior in opaque hardware under zero-trust supply-chain assumptions.
\item \textbf{ML-assisted Trojan detection:} developed golden-chip-free, ML-assisted side-channel techniques for detecting and localizing hardware Trojans on FPGAs (Journal of Hardware and Systems Security 2026; IEEE NAECON 2023).
\item \textbf{Attacks on chip-level defenses:} reverse-engineered and systematically attacked design-for-security primitives under zero-trust threat models; demonstrated authentication bypass of ring-oscillator PUFs and unlocking of dynamically obfuscated scan chains (ACM TODAES 2026).
\item \textbf{Reverse-engineering tooling:} built an automated framework that identifies and extracts security primitives from gate-level netlists, mapping a design's defensive surface for pre-silicon security review (GLSVLSI 2026).
\item \textbf{Software-side mitigation:} built HOACS, a C-library and LLVM-based transformation that keeps secrets encoded during execution on untrusted commodity processors, and evaluated the hardware-Trojan attack surface of lightweight homomorphic obfuscation for cryptographic workloads (IACR TCHES 2026).
\item \textbf{Training and mentoring:} delivered hands-on side-channel security-testing tutorials and workshops; mentored four NSF REU undergraduates across two universities; co-taught a hands-on hardware course using FPGAs and microcontrollers.
\end{rSubsection}

\end{rSection}
